% ============================================================================
% Chapter 7 — Scene Hierarchy
% ============================================================================
\chapter{Scene Hierarchy}
\label{chap:scene}

\textbf{Header:} \texttt{scene/SceneComponents.hpp}\\
\textbf{Namespace:} \texttt{Caffeine::Scene}

The scene system provides three components that describe parent--child
relationships and rendering layers.

\section{Parent}

Attaching a \texttt{Parent} component to an entity makes it a child of
another entity. The scene system reads this and computes the
\texttt{WorldTransform} for the entire hierarchy every frame.

\begin{lstlisting}
struct Parent {
    ECS::Entity parent = ECS::Entity::INVALID;
    bool        dirty  = true;   // set to true when the parent changes
};
\end{lstlisting}

\begin{lstlisting}
Entity sword = world.create("Sword");
world.add<Transform>(sword);
world.add<Parent>(sword, Parent{ playerEntity });
\end{lstlisting}

When \texttt{dirty} is true the scene system recomputes the world matrix.
You can set it manually to force a recalculation, but normally the system
manages it.

\section{WorldTransform}

Stores the precomputed world-space matrix for the entity. Written by the
scene system; read by the renderer.

\begin{lstlisting}
struct WorldTransform {
    Mat4 matrix = Mat4::identity();
};
\end{lstlisting}

Do not write to this yourself — let the scene system maintain it. Read it
when you need the final world-space transform of an entity (e.g.\ for
spawning a bullet at the gun barrel position):

\begin{lstlisting}
if (auto* wt = world.get<WorldTransform>(gunBarrel)) {
    Vec3 spawnPos = wt->matrix.transformPoint(Vec3::zero());
    // spawn bullet at spawnPos
}
\end{lstlisting}

\section{EntityLayer}

Controls the render layer. Lower layer values render behind higher values.

\begin{lstlisting}
struct EntityLayer {
    u8 layer = 0;   // 0-255
};
\end{lstlisting}

\begin{lstlisting}
world.add<EntityLayer>(background, EntityLayer{0});
world.add<EntityLayer>(player,     EntityLayer{10});
world.add<EntityLayer>(hud,        EntityLayer{100});
\end{lstlisting}

\section{Setting Up a Parent--Child Hierarchy}

\begin{lstlisting}
using namespace Caffeine::ECS;
using namespace Caffeine::Scene;

// Create parent
Entity car = world.create("Car");
world.add<Transform>(car, Transform{{200, 300, 0}});
world.add<WorldTransform>(car);

// Create child
Entity wheel = world.create("Wheel_FL");
world.add<Transform>(wheel, Transform{{-30, 20, 0}});  // local offset
world.add<WorldTransform>(wheel);
world.add<Parent>(wheel, Parent{car});
world.add<EntityLayer>(wheel, EntityLayer{5});
\end{lstlisting}

The scene system traverses the hierarchy each frame, multiplying local
transforms down the parent chain to produce each \texttt{WorldTransform}.
